\documentclass{amsart}
% For dealing with references we use the comment environment
\usepackage{verbatim}
\newenvironment{reference}{\comment}{\endcomment}
%\newenvironment{reference}{}{}
\newenvironment{slogan}{\comment}{\endcomment}
\newenvironment{history}{\comment}{\endcomment}

% For commutative diagrams we use Xy-pic
\usepackage[all]{xy}

% We use 2cell for 2-commutative diagrams.
\xyoption{2cell}
\UseAllTwocells

% We use multicol for the list of chapters between chapters
\usepackage{multicol}

% This is generally recommended for better output
\usepackage{lmodern}
\usepackage[T1]{fontenc}

% For cross-file-references

% Package for hypertext links:
\usepackage{hyperref}

% For any local file, say "hello.tex" you want to link to please

% Theorem environments.
%
\theoremstyle{plain}
\newtheorem{theorem}[subsection]{Theorem}
\newtheorem{proposition}[subsection]{Proposition}
\newtheorem{lemma}[subsection]{Lemma}

\theoremstyle{definition}
\newtheorem{definition}[subsection]{Definition}
\newtheorem{example}[subsection]{Example}
\newtheorem{exercise}[subsection]{Exercise}
\newtheorem{situation}[subsection]{Situation}

\theoremstyle{remark}
\newtheorem{remark}[subsection]{Remark}
\newtheorem{remarks}[subsection]{Remarks}

\numberwithin{equation}{subsection}

% Macros
%
\def\lim{\mathop{\mathrm{lim}}\nolimits}
\def\colim{\mathop{\mathrm{colim}}\nolimits}
\def\Spec{\mathop{\mathrm{Spec}}}
\def\Hom{\mathop{\mathrm{Hom}}\nolimits}
\def\Ext{\mathop{\mathrm{Ext}}\nolimits}
\def\SheafHom{\mathop{\mathcal{H}\!\mathit{om}}\nolimits}
\def\SheafExt{\mathop{\mathcal{E}\!\mathit{xt}}\nolimits}
\def\Sch{\mathit{Sch}}
\def\Mor{\mathop{\mathrm{Mor}}\nolimits}
\def\Ob{\mathop{\mathrm{Ob}}\nolimits}
\def\Sh{\mathop{\mathit{Sh}}\nolimits}
\def\NL{\mathop{N\!L}\nolimits}
\def\CH{\mathop{\mathrm{CH}}\nolimits}
\def\proetale{{pro\text{-}\acute{e}tale}}
\def\etale{{\acute{e}tale}}
\def\QCoh{\mathit{QCoh}}
\def\Ker{\mathop{\mathrm{Ker}}}
\def\Im{\mathop{\mathrm{Im}}}
\def\Coker{\mathop{\mathrm{Coker}}}
\def\Coim{\mathop{\mathrm{Coim}}}

% Boxtimes
%
\DeclareMathSymbol{\boxtimes}{\mathbin}{AMSa}{"02}

%
% Macros for moduli stacks/spaces
%
\def\QCohstack{\mathcal{QC}\!\mathit{oh}}
\def\Cohstack{\mathcal{C}\!\mathit{oh}}
\def\Spacesstack{\mathcal{S}\!\mathit{paces}}
\def\Quotfunctor{\mathrm{Quot}}
\def\Hilbfunctor{\mathrm{Hilb}}
\def\Curvesstack{\mathcal{C}\!\mathit{urves}}
\def\Polarizedstack{\mathcal{P}\!\mathit{olarized}}
\def\Complexesstack{\mathcal{C}\!\mathit{omplexes}}
% \Pic is the operator that assigns to X its picard group, usage \Pic(X)
% \Picardstack_{X/B} denotes the Picard stack of X over B
% \Picardfunctor_{X/B} denotes the Picard functor of X over B
\def\Pic{\mathop{\mathrm{Pic}}\nolimits}
\def\Picardstack{\mathcal{P}\!\mathit{ic}}
\def\Picardfunctor{\mathrm{Pic}}
\def\Deformationcategory{\mathcal{D}\!\mathit{ef}}

\setlength{\emergencystretch}{3em}
\begin{document}
\title[Stacks Project, Tag 0G5A]{265 --- Topological duality between coherent cohomology and compactly supported cohomology}
\maketitle
\noindent Copyright (C) 2005--2025 Johan de Jong. Stacks Project authors. Source: \href{https://stacks.math.columbia.edu/tag/0G5A}{Tag 0G5A}.

\noindent Editorial difficulty note (not author text). Difficulty H2: The proof extends a coherent complex to a compactification and represents ordinary cohomology as a colimit of boundary-ideal Ext groups. Derived bidual evaluation and proper duality convert this into an inverse-limit model for compact support. The limit/colimit passage produces the asserted topological vector-space duality..

\noindent Editorial provenance note (not author text). History: excerpt from revision \texttt{a0444\allowbreak 6e57e\allowbreak c1fbc\allowbreak 252a8\allowbreak 71afc\allowbreak ec775\allowbreak 2fb28\allowbreak 07b14}, duality.tex, lines 8898--8970. Reference presentation was adapted; original-author context and core support blocks were retained or added. The mathematical wording is unchanged. Licensed under GNU FDL 1.2 or later; no invariant sections or cover texts. See ../licenses/COPYING. Context blocks reproduce author definitions and supporting results; they are not additional counted problems.

\section{Duality for compactly supported cohomology}
\label{section-duality-compactly-supported}

\noindent
Let $k$ be a field. Let $U$ be a separated scheme of finite type over $k$.
Let $K$ be an object of $D^b_{\textit{Coh}}(\mathcal{O}_U)$. Let us define
the compactly supported cohomology $H^i_c(U, K)$ of $K$ as follows.
Choose an open immersion $j : U \to X$ into a scheme proper over $k$
and a Deligne system $(K_n)$ for $j : U \to X$ whose restriction
to $U$ is constant with value $K$. Then we set
$$
H^i_c(U, K) = \lim H^i(X, K_n)
$$
We view this as a topological $k$-vector space using the limit topology
(see More on Algebra, Section \href{https://stacks.math.columbia.edu/tag/07E7}{section topological ring (Tag 07E7)}).
There are several points to make here.

\medskip\noindent
First, this definition is independent of the choice of $X$ and $(K_n)$.
Namely, if $p : U \to \Spec(k)$ denotes the structure morphism, then
we already know that $Rp_!K = (R\Gamma(X, K_n))$
is well defined up to pro-isomorphism in $D(k)$ hence so is the limit
defining $H^i_c(U, K)$.

\medskip\noindent
Second, it may seem more natural to use the expression
$$
H^i(R\lim R\Gamma(X, K_n)) = R\Gamma(X, R\lim K_n)
$$
but this would give the same answer: since the $k$-vector spaces
$H^j(X, K_n)$ are finite dimensional, these inverse systems satisfy
Mittag-Leffler and hence $R^1\lim$ terms of Cohomology, Lemma
\href{https://stacks.math.columbia.edu/tag/0D60}{lemma RGamma commutes with Rlim (Tag 0D60)} vanish.

\medskip\noindent
If $U' \subset U$ is an open subscheme, then there is a canonical map
$$
H^i_c(U', K|_{U'}) \longrightarrow H^i_c(U, K)
$$
functorial for $K$ in $D^b_{\textit{Coh}}(\mathcal{O}_U)$.
See for example Remark \href{https://stacks.math.columbia.edu/tag/0G56}{remark covariance open lower shriek (Tag 0G56)}.
In fact, using Remark \href{https://stacks.math.columbia.edu/tag/0G57}{remark covariance etale lower shriek (Tag 0G57)}
we see that more generally such a map exists for an \'etale morphism
$U' \to U$ of separated schemes of finite type over $k$.

\medskip\noindent
If $V$ is a $k$-vector space then we put a topology on $\Hom_k(V, k)$
as follows: write $V = \bigcup V_i$ as the filtered union of its finite
dimensional $k$-subvector spaces and use the limit topology on
$\Hom_k(V, k) = \lim \Hom_k(V_i, k)$. If $\dim_k V < \infty$ then
the topology on $\Hom_k(V, k)$ is discrete. More generally, if
$V = \colim_n V_n$ is written as a directed colimit of finite dimensional
vector spaces, then $\Hom_k(V, k) = \lim \Hom_k(V_n, k)$ as topological
vector spaces.



\begin{lemma}
\label{lemma-duality-compact-support}
Let $p : U \to \Spec(k)$ be separated of finite type where $k$ is a field.
Let $\omega_{U/k}^\bullet = p^!\mathcal{O}_{\Spec(k)}$.
There are canonical isomorphisms
$$
\Hom_k(H^i(U, K), k) =
H^{-i}_c(U, R\SheafHom_{\mathcal{O}_U}(K, \omega_{U/k}^\bullet))
$$
of topological $k$-vector spaces
functorial for $K$ in $D^b_{\textit{Coh}}(\mathcal{O}_U)$.
\end{lemma}

\begin{proof}
Choose a compactification $j : U \to X$ over $k$. Let
$\mathcal{I} \subset \mathcal{O}_X$ be a quasi-coherent ideal
sheaf with $V(\mathcal{I}) = X \setminus U$.
By Derived Categories of Schemes, Proposition \href{https://stacks.math.columbia.edu/tag/0FDB}{proposition DCoh (Tag 0FDB)}
we may choose $M \in D^b_{\textit{Coh}}(\mathcal{O}_X)$
with $K = M|_U$. We have
$$
H^i(U, K) =
\Ext^i_U(\mathcal{O}_U, M|_U) =
\colim \Ext^i_X(\mathcal{I}^n, M) =
\colim H^i(X, R\SheafHom_{\mathcal{O}_X}(\mathcal{I}^n, M))
$$
by Lemma \href{https://stacks.math.columbia.edu/tag/0G2H}{lemma lift map (Tag 0G2H)}.
Since $\mathcal{I}^n$ is a coherent $\mathcal{O}_X$-module,
we have $\mathcal{I}^n$ in $D^-_{\textit{Coh}}(\mathcal{O}_X)$,
hence $R\SheafHom_{\mathcal{O}_X}(\mathcal{I}^n, M)$ is in
$D^+_{\textit{Coh}}(\mathcal{O}_X)$ by
Derived Categories of Schemes, Lemma
\href{https://stacks.math.columbia.edu/tag/0D0C}{lemma coherent internal hom (Tag 0D0C)}.

\medskip\noindent
Let $\omega_{X/k}^\bullet = q^!\mathcal{O}_{\Spec(k)}$ where
$q : X \to \Spec(k)$ is the structure morphism, see
Section \href{https://stacks.math.columbia.edu/tag/0FVU}{section duality proper over field (Tag 0FVU)}. We find that
\begin{align*}
\Hom_k(
&
H^i(X, R\SheafHom_{\mathcal{O}_X}(\mathcal{I}^n, M)), k) \\
& =
\Ext^{-i}_X(R\SheafHom_{\mathcal{O}_X}(\mathcal{I}^n, M),
\omega_{X/k}^\bullet) \\
& =
H^{-i}(X, R\SheafHom_{\mathcal{O}_X}(R\SheafHom_{\mathcal{O}_X}(
\mathcal{I}^n, M), \omega_{X/k}^\bullet))
\end{align*}
by Lemma \href{https://stacks.math.columbia.edu/tag/0FVV}{lemma duality proper over field (Tag 0FVV)}. By
Lemma \href{https://stacks.math.columbia.edu/tag/0G4I}{lemma internal hom evaluate isom (Tag 0G4I)} part (1) the canonical map
$$
R\SheafHom_{\mathcal{O}_X}(M, \omega_{X/k}^\bullet)
\otimes_{\mathcal{O}_X}^\mathbf{L} \mathcal{I}^n
\longrightarrow
R\SheafHom_{\mathcal{O}_X}(R\SheafHom_{\mathcal{O}_X}(
\mathcal{I}^n, M), \omega_{X/k}^\bullet)
$$
is an isomorphism. Observe that
$\omega^\bullet_{U/k} = \omega^\bullet_{X/k}|_U$
because $p^!$ is constructed as $q^!$ composed with restriction to $U$.
Hence $R\SheafHom_{\mathcal{O}_X}(M, \omega_{X/k}^\bullet)$ is
an object of $D^b_{\textit{Coh}}(\mathcal{O}_X)$ which restricts
to $R\SheafHom_{\mathcal{O}_U}(K, \omega_{U/k}^\bullet)$ on $U$.
Hence by Lemma \href{https://stacks.math.columbia.edu/tag/0G4U}{lemma tensoring Deligne system (Tag 0G4U)} we conclude that
$$
\lim H^{-i}(X, R\SheafHom_{\mathcal{O}_X}(M, \omega_{X/k}^\bullet)
\otimes_{\mathcal{O}_X}^\mathbf{L} \mathcal{I}^n)
$$
is an avatar for the right hand side of the equality of the lemma.
Combining all the isomorphisms obtained in this manner we get
the isomorphism of the lemma.
\end{proof}
\end{document}
